% !TeX root = ../抽象代数笔记.tex
% 由主文档通过 \input 载入；请编译主文档。
\setcounter{chapter}{4} % 保留原书章号。
\chapter{正规子群、商群与同态定理}\label{ch:quotient-groups}

\begin{description}
  \item[学习目标] 理解正规性控制商群乘法的良定义性，能够用核与像识别商群。
  \item[先修知识] 第\ref{ch:cosets}章。
\end{description}

\section{群同态、核与像}
在定义\ref{def:group-isomorphism}的基础上系统引入群同态，证明同态保持单位元、逆元和整数次幂。对同态 $f:G\to H$ 定义核 $\ker f$ 与像 $\operatorname{im}f$，证明单射等价于核只含单位元。结合循环群、符号映射和行列式映射计算核与像。

\section{正规子群与商群构造}
定义正规子群（normal subgroup）$N\trianglelefteq G$，证明共轭不变与左右陪集相等的等价性。证明陪集乘法 $(gN)(hN)=ghN$ 良定义当且仅当 $N$ 正规，再定义商群（quotient group）$G/N$。说明正规性不同于交换性，并证明任意同态的核正规。

\section{商映射的分解性质与第一同构定理}
对正规子群 $N\trianglelefteq G$，定义自然投影 $\pi:G\to G/N$。若同态 $f:G\to H$ 满足 $N\subseteq\ker f$，证明存在唯一同态 $\overline f:G/N\to H$ 使 $f=\overline f\circ\pi$。由此证明第一同构定理 $G/\ker f\cong\operatorname{im}f$，并强调需构造映射、验证良定义，再证明双射。

\section{其余同构定理与对应定理}
对 $H\leq G$、$N\trianglelefteq G$，证明 $HN/N\cong H/(H\cap N)$。对 $N\subseteq K$ 且 $N,K\trianglelefteq G$，证明 $(G/N)/(K/N)\cong G/K$。建立 $G/N$ 的子群与 $G$ 中包含 $N$ 的子群之间的对应，同时说明正规性如何传递。

\section{例题、练习与自检}
计算 $\mathbb Z/m\mathbb Z$、$S_n/A_n$（$n\geq2$）及若干矩阵群同态的像。证明指数为 $2$ 的子群必正规，完成 $A_4$ 不存在 $6$ 阶子群的反例。安排一个非正规子群的陪集乘法失败算例，使良定义性成为实际检查步骤。

